\section{Accessibility of locally constant coefficient theories}
\providecommand{\sSet}{\mathrm{sSet}}
\providecommand{\cC}{\mathcal C}
\providecommand{\cA}{\mathcal A}
\providecommand{\cU}{\mathcal U}
\providecommand{\cV}{\mathcal V}
\providecommand{\D}{\mathbb D}
\providecommand{\PA}{P_{\mathcal A}}

\begin{claim}{H.1}{Proved}{Products of derivators and intersections of localizers}
Let $I\in\cV$ be a set and let $\D_i$ be derivators on $\cU$-small categories with values in $\cV$-categories. Their pointwise product is a derivator, and
\[
 W^{\mathrm{lc}}_{\prod_{i\in I}\D_i}
 =\bigcap_{i\in I}W^{\mathrm{lc}}_{\D_i}.
\]
This includes the empty product.
\end{claim}
\begin{proof}
All restrictions and Kan extensions are taken coordinatewise. The coproduct-to-product comparison of Der~1, the detection of isomorphisms of Der~2, the two adjunctions of Der~3, and the pointwise mate isomorphisms of Der~4 are therefore products of their corresponding comparisons in the $\D_i$. A product of isomorphisms is an isomorphism, and all the indicated products exist in the larger universe. This verifies the four axioms without a strongness assumption.

An object of the product derivator is locally constant precisely when every coordinate is locally constant, so its locally constant category is the product of the locally constant categories. Each such category is nonempty: the terminal object, constructed by right Kan extension from the empty category, is constant and hence locally constant. For functors $F_i:C_i\to E_i$ between nonempty categories, $\prod F_i$ is an equivalence if and only if all $F_i$ are equivalences. To recover full faithfulness of one factor, fill the other coordinates by chosen objects and their identity morphisms; injectivity and surjectivity of the resulting hom-set map give the two required properties. To recover essential surjectivity, fill the other target coordinates with objects already in the corresponding images. This proves the asserted equality. The empty product is the terminal derivator, whose localizer is all functors, as is the empty intersection.
\end{proof}

\begin{claim}{H.2}{Proved}{Vop\v{e}nka's principle gives accessibility in the indexing universe}
Assume Vop\v{e}nka's principle for $\cU$-classes: every $\cV$-class of $\cU$-small structures which is proper relative to $\cU$ satisfies the principle. Then $W^{\mathrm{lc}}_\D$ is generated as a fundamental localizer by a $\cU$-small set, for every derivator in the stated two-universe setting.
\end{claim}
\begin{proof}
We check the arbitrary-derivator input of GALOIS B.3. For a small category $T$ and a subcategory $R$, the full subprederivator of $\D^T$ inverting the $R$-arrows has both adjoints to its inclusion. Cisinski's Theorem 3.6 identifies it with cocontinuous morphisms out of the simplicial-presheaf derivator localized at $R$. The inclusion of the local simplicial presheaves preserves homotopy colimits, because these are pointwise and a pointwise colimit still inverts $R$. Precomposition with the two cocontinuous adjoints gives the left projector; the opposite derivator gives the right projector, and the same argument on every shift makes them morphisms of derivators. These are the steps of \cite[Theorem 3.6, Lemma 3.9, and Theorem 3.13]{FGHCisinski2009} actually required here.

For a Grothendieck opfibration $T\to S$, apply the right projector to the vertical arrows. Restriction to a fibre commutes with this projector: the left adjoint of fibre restriction takes a locally constant object to a vertically locally constant object, since its value at $x\in T_t$ is the coproduct of the values on the transports of $x$, indexed by arrows $t\to s$. Taking right adjoints gives the claimed commutation. For a transport-preserving map of opfibrations, right Kan extension commutes with fibre restriction; the inclusion of the strict-fibre comma category in the total comma category has the left adjoint given by cocartesian transport. The unit and counit of restriction followed by projected right Kan extension can consequently be tested fibrewise. If all fibre restrictions are equivalences, Der~2 makes the global unit and counit isomorphisms. Equivalence on the vertically constant categories restricts to equivalence on the fully locally constant categories, since the transport maps are reflected as isomorphisms by the fibre equivalences. Finally any $A\to S$ has the canonical replacement $\int_{s\in S}(A/s)\to S$; its forgetful functor to $A$ has an adjoint section and induces equivalence on locally constant objects. This proves relative Theorem A with precisely the Der~1--4 hypotheses in the question. Saturation is immediate from equivalence of the specified restriction functors, and a terminal object gives the usual adjoint contraction. Thus $W^{\mathrm{lc}}_\D$ is a fundamental localizer.

Cisinski's Scholie 4.2.17 states that, under Vop\v{e}nka, every fundamental localizer is accessible. Its proof uses Theorem 4.2.15 to pass to the simplicial-presheaf localizer for the test category $\Delta$, Scholie 1.4.6 to obtain a small generating family there, and the accessibility-preserving inverse correspondence in Theorem 4.2.15 to return to $\Cat$. Apply these three statements inside the indexing universe $\cU$. The resulting generators are $\cU$-small, despite the coefficient categories living in $\cV$. No bound on the size of the family of coefficient objects is inserted. The explicitly stated class form of Vop\v{e}nka is relevant: an internal first-order scheme with no permission to use the externally supplied derivator as a class parameter would not by itself justify this application.
\end{proof}

\begin{claim}{H.3}{Proved}{A nonaccessible locally constant localizer when there are no measurable cardinals}
Suppose that $\cU$ contains no measurable cardinal. There exists a derivator on all $\cU$-small categories with values in $\cV$-categories for which $W^{\mathrm{lc}}$ is not generated by a $\cU$-small set. It can be chosen as the derivator of a complete and cocomplete, locally $\cU$-small infinity-category of localized $1$-types.
\end{claim}
\begin{proof}
For each infinite cardinal $\kappa\in\cU$, put
\[
 A_\kappa=\ZZ^\kappa/\ZZ^{<\kappa},
\]
where the denominator consists of functions supported on fewer than $\kappa$ coordinates, and let $\cA$ be the class of these groups. A group $G$ is $\cA$-reduced when every homomorphism $A_\kappa\to G$ is trivial. We use exactly Proposition 6.2 of \cite{FGHCSS2005}: in the absence of measurable cardinals, $\ZZ$ is $\cA$-reduced. The statement relativizes to $\cU$, since a nonzero homomorphism $A_\kappa\to\ZZ$ would yield a measurable cardinal at most $\kappa$. The paper's Theorem 6.4 concerns a homotopy idempotent functor; the additional derivator construction and the shift from $\pi_1$ to $\pi_2$ below are necessary here. The reflector itself is also covered by \cite[Theorem 3.3 with $n=1$, and Section 6.1]{FGHChristensen2024}; that theorem does not by itself compute the locally constant coefficient localizer used below.

We first construct the coefficient category with its size bounds. For a $\cU$-small group $G$, let $R(G)$ be the intersection of the kernels of all homomorphisms from $G$ to $\cA$-reduced groups. One may use only quotient images of $G$: reduced groups are closed under subgroups, and all such quotient images have size at most $|G|$. There is a $\cU$-small set of these quotients. The group $G/R(G)$ embeds into their product, so it is reduced, since products and subgroups of reduced groups are reduced. This constructs the reflector $\PA G=G/R(G)$ within $\cU$. Its universal property is immediate from the definition of $R(G)$.

The assignment is functorial: for a homomorphism $v:G\to H$, composing every map from $H$ to a reduced group with $v$ gives $v(R(G))\subseteq R(H)$. Apply the quotient at every isotropy group of a groupoid. Conjugation along an arrow $x\to y$ identifies $R(\operatorname{Aut}(x))$ with $R(\operatorname{Aut}(y))$. Thus these subgroups form a normal subgroupoid; quotient each arrow torsor by its source subgroup, equivalently by its target subgroup. Composition is well defined, and the object set is unchanged. For a reduced groupoid $H$, every functor $G\to H$ factors uniquely through this quotient, and the same is true of natural isomorphisms. Thus it defines a reflector of infinity-categories of $1$-types. Compose it with $1$-truncation to obtain
\[
 L:\mathcal S_{\cU}\longrightarrow\cC,
 \qquad LX=N\PA\Pi_1X,
\]
where $\cC$ is the full infinity-category of $\cU$-small $1$-types with reduced isotropy groups. Its object collection and all its diagram homotopy categories are $\cV$-small; its mapping spaces are $\cU$-small.

The reflector preserves products. For groups, a homomorphism $G\times H\to K$ into a reduced group descends on both commuting factors to $\PA G$ and $\PA H$, and their descended images still commute. The product $\PA G\times\PA H$ is itself reduced, proving the assertion by the universal property. The groupoid and space statements follow componentwise and from product preservation of $\Pi_1$. The category $\cC$ is closed under all $\cU$-small homotopy limits: it is the orthogonality class of $S^2\to*$ and all $BA_\kappa\to*$. Its colimits are obtained by applying $L$ to space colimits. Therefore it is complete and cocomplete for the required shapes, and
\[
 \D_{\cC}(I)=\Ho\Fun(NI^{\mathrm{op}},\cC)
\]
satisfies Der~1--4. More explicitly, functors on a disjoint union are tuples of functors, an equivalence of natural transformations is detected at each object, both Kan extensions exist because the required comma diagrams are $\cU$-small, and their pointwise formulas are precisely Der~4. All these Kan extensions have values in $\cC$ by its established completeness and cocompleteness. Presentability in $\cU$ is not assumed.

We compute its locally constant localizer. On a connected pointed space $(X,x)$, evaluation $U_x:\Fun(X,\cC)\to\cC$ is conservative, since paths identify all point evaluations, and it preserves colimits, which are computed pointwise. It has a left adjoint. The infinity-categorical Barr--Beck theorem, \cite[Theorem 4.7.3.5]{FGHLurieHA}, says that a right adjoint is monadic if it is conservative and preserves geometric realizations of the simplicial objects split by it. Here all geometric realizations are preserved. The free-evaluation monad is
\[
 T_X(M)=(\Omega_xX)\otimes M
        =L(\Omega_xX\times M)
        =L\Omega_xX\times M.
\]
All equalities refer to the canonical maps, and the multiplication is induced by concatenation of loops. For a pointed map $f:X\to Y$, restriction is consequently restriction of modules for the specified monad map $T_X\to T_Y$. It is an equivalence if and only if that monad map is an equivalence. Necessity also follows directly by identifying the left adjoints of the two evaluation functors; evaluation of their comparison at $M=*$ gives $L\Omega f$. Sufficiency is the equivalence of module categories of equivalent monads.

This calculation concerns infinity-categories, whereas the question uses ordinary categories. The distinction causes no loss here. Restriction preserves all cotensors. Full faithfulness on homotopy categories, tested against the cotensors $N^K$ for every $\cU$-small space $K$, gives bijections $[K,\Map(M,N)]$ for every $K$, so Yoneda makes the specified mapping-space comparison a weak equivalence. Essential surjectivity is unchanged on passage to the homotopy category. Thus the two equivalence tests agree. For disconnected spaces, the initial and terminal objects of $\cC$ distinguish missing and repeated components; equivalence forces a bijection on components. The connected calculation then applies to each factor. We have proved
\[
 f\in W^{\mathrm{lc}}_{\D_{\cC}}\quad\Longleftrightarrow\quad
 \begin{cases}
 \pi_0f\text{ is bijective},\\
 \pi_1f\text{ is an isomorphism at every basepoint},\\
 \PA(\pi_2f)\text{ is an isomorphism at every basepoint}.
 \end{cases}                                                     \tag{H.3.1}
\]
There is no condition on higher homotopy groups. In particular $K(A_\kappa,2)\to*$ is in this class for every $\kappa$, since the quotient unit $A_\kappa\to\PA A_\kappa$ is itself a homomorphism to a reduced group and is therefore trivial.

Every map in (H.3.1) induces integral cohomology isomorphisms in degrees zero, one, and two. The first two degrees follow from $\pi_0$ and $\pi_1$. To see the third on a connected component, filter the bar resolution of $\pi=\pi_1X$ tensored with the cellular chains of the universal cover. The first-quadrant cohomology spectral sequence has entries $H^p(\pi;H^q(\widetilde X;\ZZ))$. Since the universal cover is simply connected, its row $q=1$ vanishes and its row $q=2$ is $\Hom(\pi_2X,\ZZ)$. The terms of total degree two therefore give the natural exact sequence
\[
 0\longrightarrow H^2(\pi;\ZZ)\longrightarrow H^2(X;\ZZ)
 \longrightarrow\Hom_\pi(\pi_2X,\ZZ)
 \longrightarrow H^3(\pi;\ZZ).
\]
Because $\ZZ$ is reduced, reflection gives $\Hom(\PA B,\ZZ)=\Hom(B,\ZZ)$ naturally in $B$. Thus (H.3.1) makes all terms other than $H^2(X;\ZZ)$ isomorphic in the comparison of these exact sequences, including the group actions; a diagram chase gives the remaining isomorphism. Disconnected cohomology is the product of these component calculations.

Suppose now that a $\cU$-small set $S$ of functors generates $W^{\mathrm{lc}}_{\D_{\cC}}$. Choose a regular uncountable $\kappa\in\cU$ larger than the total number of simplices in the nerves of the sources and targets of every member of $S$. For an abelian group $B$, write $R_\kappa(B)=B^\kappa/B^{<\kappa}$. This functor is exact: products of surjections are surjective, and a lift of a vector supported on fewer than $\kappa$ coordinates can be chosen with the same support. Applying the snake lemma to the inclusion of the bounded-support sequence in the product sequence gives exactness of the quotient sequence.

If $C_*$ is a free chain complex with fewer than $\kappa$ basis elements in total, there is a natural degreewise isomorphism
\[
 \Hom(C_*,A_\kappa)
 \cong R_\kappa\bigl(\Hom(C_*,\ZZ)\bigr).
\]
Surjectivity follows by lifting each basis value to $\ZZ^\kappa$. For injectivity, fewer than $\kappa$ basis values have supports of size less than $\kappa$; their union still has size less than $\kappa$ by regularity. The isomorphism commutes with the differential. Exactness of $R_\kappa$ then gives
\[
 H^n(C_*;A_\kappa)\cong R_\kappa H^n(C_*;\ZZ).
\]
Apply this to the normalized nerve-chain complexes of all members of $S$. Their integral cohomology comparisons in degrees zero, one, and two are isomorphisms, so the same is true with coefficients $A_\kappa$. It follows that $Z=K(A_\kappa,2)$ is local for every member of $S$: the mapping spaces into this Eilenberg--Mac Lane space are $2$-types with homotopy groups $H^{2-i}(-;A_\kappa)$ for $0\leq i\leq2$; all components are translates of one another.

For any fixed space $Z$, the class of functors inducing equivalences on mapping spaces into $Z$ is a fundamental localizer. To check relative Theorem A directly, use the natural decomposition $BA\simeq\hocolim_sB(A/s)$ and the analogous decomposition for $B$; mapping into $Z$ changes this homotopy colimit to a homotopy limit. Its other two axioms follow immediately. Consequently a space local for $S$ inverts the whole fundamental localizer generated by $S$. In particular it inverts $K(A_\kappa,2)\to*$. This is impossible: the induced map on components of mapping spaces would identify the identity of $K(A_\kappa,2)$ with its constant map, whereas $A_\kappa\neq0$, as the constant function $1$ has support $\kappa$. This contradiction proves nonaccessibility.
\end{proof}

\begin{claim}{H.4}{Proved}{Represented derivators give exactly the first two truncation localizers}
Let $C$ be complete and cocomplete for $\cU$-small diagrams. For its represented derivator, $W^{\mathrm{lc}}=W_1$ when $C$ is not a preorder up to equivalence; $W^{\mathrm{lc}}=W_0$ when $C$ is a preorder with more than one isomorphism class; and $W^{\mathrm{lc}}$ is all functors when $C$ is equivalent to the terminal category.
\end{claim}
\begin{proof}
Locally constant $C$-valued diagrams on $X$ factor uniquely through the fundamental groupoid $\Pi_1X$; natural transformations factor as well. Thus their category is $\Fun(\Pi_1X,C)$, with the specified restriction functor.

Suppose first that there are distinct arrows $r,s:a\to c$ in $C$. For every nonempty set $T$, distinct projections $c^T\to c$ are distinct: prescribing $r$ and $s$ in the corresponding two coordinates gives a map from $a$ separating them. For a left $G$-set $T$, define $\rho_g:c^T\to c^T$ by $p_t\rho_g=p_{g^{-1}t}$ for every product projection $p_t$. These maps form a $G$-action. Taking the regular left $G$-set $T=G$ gives a faithful action, since a nonidentity translation moves a coordinate and distinct projections are distinct.

Consider a homomorphism $\varphi:G\to H$. If restriction $\Fun(BH,C)\to\Fun(BG,C)$ is essentially surjective, every $G$-action is isomorphic to a restricted $H$-action. Its kernel therefore acts trivially on the faithful permutation action, so $\varphi$ is injective. Identify $G$ with its image in $H$. If it is a proper subgroup, let $H$ act on $c^{H/G}$ by permutation of the left cosets, and give $c$ the trivial action. Projection onto the coordinate $G\in H/G$ is $G$-equivariant. It is not $H$-equivariant, since some element of $H$ moves that coset and distinct projections are distinct. Hence restriction is not full. Equivalence of the action categories therefore forces $\varphi$ to be an isomorphism; the converse is immediate.

For a general functor of groupoids, failure of surjectivity on components is detected by two natural transformations that differ by $r$ and $s$ on a missed component. Failure of injectivity on components is detected by a diagram taking nonisomorphic objects of $C$ on two components that merge. Such objects exist: if every object were isomorphic to a terminal object, every hom-set would be a singleton, contrary to the chosen $r,s$. Thus equivalence of the restriction categories first forces a bijection of components and then reduces to the group calculation. This proves $W^{\mathrm{lc}}=W_1$.

If $C$ is a preorder, replace it by its equivalent poset. Every group action is trivial, so the local-system category is $C^{\pi_0X}$. When this poset is nontrivial, its initial and terminal elements are distinct. For a noninjective component function, a tuple taking these two different values on a pair of merged components cannot lie in the essential image of restriction. For a nonsurjective component function, choose two target tuples which agree on its image and take terminal and initial values, in that order, on a missed component. There is no arrow between the target tuples in that direction, while their restrictions are equal. Hence restriction is not fully faithful. Conversely a bijection of components gives an isomorphism of product categories. The terminal-category case is immediate.
\end{proof}
